\section{A polar-retaining completion and the actual positive source column}
\label{sec:polar-source-completion}

The narrow inherited intersection in MA-MO-007 loses both source
polar modes. We now retain those modes by completing actual source
sections together with the inherited diagonal norm. This produces
a bounded representation and an exact quotient map. The original
smoothed positive source column extends to a trace-class map into
the completed common space. The target of that map is kept explicit:
it is not automatically an endomorphism of the completed positive
source image.

Keep the rational place set \(S=\{\infty\}\), the centered theta
map (MA-MO.83), and the original smooth common-sum space
\(\mathscr S_{\rm sum}\) in (MA-MO.78). Put
\[
 \begin{aligned}
 H&=H_{\Sigma,a,s},&0<a&<1/2,&s&\ge0,\\
 I&=\Theta\mathscr U_0,&J&=\overline I^{\,H},
       &K&=K_{a,s,\phi},\\
 \mathscr H_+&=i_+\mathscr S_{\rm even}(\mathbb R),
       &\mathscr H_-&=i_-\mathcal E.
 \end{aligned}
 \tag{MA-MO.250}
\]
Thus \(i_+f=(\Theta f,R\Theta\mathfrak Ff)\),
\(i_-u=(u,u)\), and \(\mathscr U_0\) imposes
\(f(0)=\int f=0\). Let
\[
 \begin{aligned}
 f_0(t)&=(1-2\pi t^2)e^{-\pi t^2},\\
 f_1(t)&=2\pi t^2e^{-\pi t^2},&s_j&=i_+f_j.
 \end{aligned}
 \tag{MA-MO.251}
\]
These are actual even Schwartz sources, with
\((f_j(0),\int f_j)=e_j\). No bare polar exponential is
being declared an element of a narrow weighted \(L^2\) space.

\paragraph{MA-MO-016 (equivariant polar completion and actual source-column
transport).}
The map
\[
 \begin{aligned}
 \mathfrak s(u,c)&=i_-u+c_0s_0+c_1s_1,\\
 \mathcal W&=H\oplus\mathbb C^2,&
 \mathcal W_+&=J\oplus\mathbb C^2,&
 \mathcal W_-&=H\oplus\{0\}
 \end{aligned}
 \tag{MA-MO.252}
\]
identifies \(\mathcal W\) with a Hilbert completion of
\(\mathscr S_{\rm sum}\), and \(\mathcal W_+\) and
\(\mathcal W_-\) with its closed positive and negative source
image completions. The embedding \(\mathfrak s\) into pairs
of locally square-integrable functions is continuous and injective.
The diagonal norm is precisely the inherited norm \(H\).

Simultaneous translations of the actual pair functions extend
to a strongly continuous representation on \(\mathcal W\).
Both \(\mathcal W_+\) and \(\mathcal W_-\) are invariant.
There are equivariant Hilbert exact sequences
\[
 \begin{gathered}
 0\longrightarrow J\longrightarrow\mathcal W_+
       \xrightarrow{\,\operatorname{pol}\,}\mathbb C^2
       \longrightarrow0,\\
 0\longrightarrow\mathcal W_+\longrightarrow\mathcal W
       \xrightarrow{\,\mathfrak q\,}H/J\longrightarrow0,
       \qquad \mathfrak q(u,c)=u+J.
 \end{gathered}
 \tag{MA-MO.253}
\]
The last quotient identification is isometric in the displayed
direct-sum norm, and its representation is the actual quotient
translation representation. Also
\(\mathcal W_+\cap\mathcal W_-=J\) and
\(\mathcal W_++\mathcal W_-=\mathcal W\).
Changing the two actual source sections gives a boundedly
isomorphic completion, by an isomorphism equal to the identity
on the original common-sum space.

Assume either analytic regime of MA-MO-010, with its stated test
class, and write \(D=P_+-P_-\) for the original source matrix.
Then the actual smooth map \(C_gD|_{\mathscr H_+}\) extends
uniquely to
\[
       \mathcal B_g:\mathcal W_+\longrightarrow\mathcal W,
                      \qquad \mathcal B_g\in\mathcal S_1.
 \tag{MA-MO.254}
\]
Its restriction to \(J\) takes values in \(\mathcal W_+\).
The possible outgoing quotient defect has rank at most two.
More precisely, define actual smooth functions
\[
       r_j=(1-\phi)(s_j)_0+\phi(s_j)_1\in\mathcal E.
 \tag{MA-MO.255}
\]
Then
\[
       \mathfrak q\mathcal B_g(y,c)
             =-\sum_{j=0}^1c_j\,[C_gr_j]\quad\hbox{in }H/J.
 \tag{MA-MO.256}
\]
In particular \(\mathcal B_g\) preserves \(\mathcal W_+\)
if and only if both \(C_gr_j\) belong to \(J\). The theorem
does not assert that condition. Nor does it extend the entire
original matrix to \(\mathcal W\), assign an endomorphism
trace to (MA-MO.254), or identify a quotient trace with a local
explicit formula.

\begin{proof}
\emph{Actual sections and the exact algebraic decomposition.}
The Gaussian integral and its derivative give
\(\int e^{-\pi t^2}dt=1\) and
\(\int2\pi t^2e^{-\pi t^2}dt=1\), proving the stated
moment normalization in (MA-MO.251). In fact
\(\mathfrak Ff_0=f_1\) and \(\mathfrak Ff_1=f_0\),
though only the moment normalization is needed below.
Poisson summation (MA-MO.84) gives
\[
 (i_+f)_0-(i_+f)_1
      =-f(0)e^{x/2}+\left(\int f\right)e^{-x/2}.
 \tag{MA-MO.257}
\]
Use the two coefficients \((f(0),\int f)\) as the polar
coordinates. The corresponding two exponentials are linearly
independent in distributions. Every pair in
\(\mathscr S_{\rm sum}\) has unique such coordinates;
subtracting the matching combination of \(s_0,s_1\)
makes its components equal. The equal component lies in
\(\mathcal S_>\cap\mathcal S_<=\mathcal E\), as proved
in MA-MO-007. Conversely every pair in (MA-MO.252) with
\(u\in\mathcal E\) belongs to the original common-sum
space. Thus
\[
 \mathscr S_{\rm sum}=\mathfrak s(\mathcal E\oplus\mathbb C^2),
 \qquad
 \mathscr H_+=\mathfrak s(I\oplus\mathbb C^2),
 \qquad
 \mathscr H_-=\mathfrak s(\mathcal E\oplus\{0\}).
 \tag{MA-MO.258}
\]
For the middle equality, write any even Schwartz source as
\(f=h+f(0)f_0+(\int f)f_1\), with \(h\in\mathscr U_0\).
Then \(i_+h=i_-\Theta h\). This also proves the reverse
inclusion and uses no hypothesized theta surjectivity beyond
the definition of \(I\).

On each compact logarithmic interval the weight of \(H\)
is bounded below by a positive constant. Thus the diagonal
\(H\) component maps continuously to local \(L^2\), and
the fixed sections do also. Injectivity of \(\mathfrak s\)
follows first from (MA-MO.257), which forces both coefficients
to vanish, and then from its diagonal component. This proves
the Hilbert realization and the asserted closures in
(MA-MO.252). If other normalized actual source sections are
chosen, their differences from \(s_j\) are diagonal vectors
in \(I\); the coordinate change is
\((u,c)\mapsto(u+Tc,c)\), with \(T:\mathbb C^2\to I\)
bounded. Its inverse has the opposite finite column.
Unnormalized independent moment sections differ additionally
by an invertible two-by-two matrix. This proves the claimed
section independence up to the specified bounded isomorphism.

\emph{Translations and their actual theta cocycle.}
Set
\[
 \begin{aligned}
 U_\tau u(x)&=u(x-\tau),&
 \rho_\tau f(t)&=e^{-\tau/2}f(e^{-\tau}t),\\
 d_0(\tau)&=e^{-\tau/2},&d_1(\tau)&=e^{\tau/2},\\
 k_j(\tau)&=\Theta\bigl(\rho_\tau f_j-d_j(\tau)f_j\bigr).
 \end{aligned}
 \tag{MA-MO.259}
\]
Direct substitution in the theta sum gives
\(\Theta\rho_\tau=U_\tau\Theta\), and Fourier scaling
gives \(\mathfrak F\rho_\tau=\rho_{-\tau}\mathfrak F\).
Hence \(i_+\rho_\tau=U_\tau i_+\), with simultaneous
translation on the pair. The two source moments of
\(\rho_\tau f\) scale by \(d_0(\tau),d_1(\tau)\).
Consequently every \(k_j(\tau)\) lies in \(I\), and the
actual coordinate formula is
\[
 \pi_\tau(u,c)=
 \left(U_\tau u+\sum_{j=0}^1c_jk_j(\tau),
                 (d_0(\tau)c_0,d_1(\tau)c_1)\right).
 \tag{MA-MO.260}
\]
This is a representation because it agrees with the actual
translations on a dense source domain; alternatively its
cocycle identity follows directly from (MA-MO.259).

For completeness the required continuity is in the declared
norms. The weight \(w\) of \(H\) satisfies
\[
       \frac{w(x+\tau)}{w(x)}
       \le e^{2a|\tau|}2^{s/2}(1+\tau^2)^{s/2}.
 \tag{MA-MO.261}
\]
Thus \(U_\tau\) is bounded, locally uniformly in \(\tau\),
and is strongly continuous by approximation with compact
smooth functions. The map \(\Theta:\mathscr U_0\to\mathcal E\)
is continuous: at the right end, estimates for the sum and
each logarithmic derivative follow from finitely many
Schwartz seminorms of the source; at the left end use
\(\Theta f=R\Theta\mathfrak Ff\) and the same estimates
for \(\mathfrak Ff\). Arbitrarily large additive decay
powers give every prescribed exponential logarithmic weight.
Since dilation is smooth on Schwartz space, \(k_j\) is
continuous in \(\mathcal E\), hence in \(H\). Formula
(MA-MO.260) therefore defines a bounded strongly continuous
representation. It preserves \(J\) by the density of its
translation-invariant subspace \(I\), and the cocycle lies
in \(J\). This proves invariance of both completed images.

The kernel and surjectivity assertions in (MA-MO.253) now
follow from the direct-sum coordinates. In the second quotient
the \(\mathbb C^2\) component may be removed at no remaining
cost, and the infimum over the first component is precisely
the quotient norm of \(H/J\). The cocycle dies in that
quotient, so its action is \(u+J\mapsto U_\tau u+J\).
The polar quotient in the first sequence has the two
characters \(d_0,d_1\). None of these statements asserts
that the displayed source-section splitting is equivariant.

\emph{The exact negative projection and finite polar columns.}
The original definition of \(p_-\) gives, in these coordinates,
\[
       p_-\mathfrak s(u,c)=u+\sum_jc_jr_j.
 \tag{MA-MO.262}
\]
Each \(r_j\) belongs to \(\mathcal E\): sufficiently far
left only \((s_j)_1\) remains, and sufficiently far right
only \((s_j)_0\) remains, with the rapid weighted tails
specified by \(\mathcal S_<\) and \(\mathcal S_>\).
The intervening region is compact and smooth. Consequently
\(p_-:\mathcal W\to H\) and
\(P_-=i_-p_-:\mathcal W\to\mathcal W_-\) extend
boundedly, and \(P_-\) remains a projection since
\(p_-i_-=1_H\). This does not assert that \(p_+\) has
a comparable bounded extension.

Compact convolution \(C_g=\int g(\tau)\pi_\tau\,d\tau\)
is bounded on \(\mathcal W\), by the local bounds just
proved. On source vectors it equals the actual simultaneous
convolution, and
\[
 C_gi_+f=i_+\mathcal D_gf,\qquad
 \mathcal D_gf(t)=\int g(\tau)e^{-\tau/2}f(e^{-\tau}t)\,d\tau.
 \tag{MA-MO.263}
\]
This is an even Schwartz function. Thus convolution preserves
\(\mathscr H_+\), \(\mathscr H_-\), and their completions.
Both \(P_+s_j\) and \(P_-s_j\) are genuine original source
outputs: the first is \(i_+p_+s_j\) with even Schwartz
\(p_+s_j\), and the second is \(i_-r_j\).
Therefore \(C_gDs_j\in\mathscr S_{\rm sum}\subset\mathcal W\)
without any narrow multiplier continuation.

\emph{The actual trace-class map on the ideal column.}
The bounded functionals of MA-MO-011 give
\(\Lambda=(\lambda_0,\lambda_1):J\to\mathbb C^2\),
with \(K=\ker\Lambda\). Choose fixed
\(h_0,h_1\in\mathscr U_0\), put \(y_j=\Theta h_j\),
and normalize \(\Lambda y_j=e_j\), as in (MA-MO.147).
The bounded map
\[
          R_\phi y=y-\sum_j\lambda_j(y)y_j:J\to K
 \tag{MA-MO.264}
\]
preserves the actual smooth domains:
\(R_\phi I=\Theta\mathscr U_{\rm pf}\).
For \(f\in\mathscr U_0\), the defining source equations
(MA-MO.85) give \(D i_+f=i_+Lf\), exactly as in
MA-MO-008. Hence the two fixed vectors
\(b_j=C_gD i_-y_j=i_+\mathcal D_gLh_j\) belong to
\(\mathscr H_+\subset\mathcal W_+\).
Let \(B_{\phi,g}:K\to J\) be the actual map constructed
in MA-MO-010. The formula
\[
 \mathcal B_g(y,c)=
   i_-B_{\phi,g}R_\phi y
       +\sum_j\lambda_j(y)b_j
       +\sum_jc_j C_gDs_j
 \tag{MA-MO.265}
\]
now defines a trace-class map \(\mathcal W_+\to\mathcal W\).
Indeed the first term is a trace-class map composed with
bounded maps; the other four columns have finite rank.
It agrees with the original operator on
\(I\oplus\mathbb C^2\) by the exact smooth decomposition
above, so density proves uniqueness. This also proves
independence from the auxiliary cutoff-moment correctors.
The trace-norm continuity in smooth fixed-support test spaces,
or boundedness for common Gevrey test bounds, follows from
MA-MO-010 for the first term and ordinary Schwartz continuity
of the finitely many actual source columns for the others.

The first two terms in (MA-MO.265) lie in \(\mathcal W_+\).
On a polar section \(\mathfrak q C_gP_+s_j=0\), since
\(P_+s_j\in\mathscr H_+\) and convolution preserves that
space. The negative term is
\(\mathfrak q C_gP_-s_j=[C_gr_j]\). This proves the
exact formula (MA-MO.256) and its stated if-and-only-if
criterion. It also explains why an endomorphism target
cannot be obtained merely by completing the two additional
source vectors.
\end{proof}

\paragraph{Source scope and the remaining matrix problem.}
The completion and equivariant exact sequences use only Poisson
summation and the declared weights; their construction has no
zero-location premise. The analytic regimes enter precisely in
the bounded moment functionals and the trace-class ideal-column
estimate inherited from MA-MO-010--011. In particular the result
does not assert unconditional smooth-test transport at every
narrow width. It changes the completion of the actual common-sum
source by retaining its polar coordinates, while retaining the
exact inherited norm on its diagonal part.

The original smooth common-space matrix is still defined on
\(\mathfrak s(\mathcal E\oplus\mathbb C^2)\). The theorem
does not provide its bounded action on the complementary
off-ideal diagonal directions in \(H\). That is the remaining
negative-column problem. A trace of a completed full matrix,
a relative trace, or a quotient trace would require those
additional maps and the corresponding trace-class comparisons.
The finite outgoing formula (MA-MO.256) is a separate condition
even on the positive source column.

The exact primary locators are Meyer (2005),
\texttt{sec:summation}, \texttt{lem:Hilm\_p\_embed}, the
definitions of \(p_-,P_-\), \texttt{eq:Poisson\_summation},
and \texttt{lem:Poisson\_sumpm}; author-manuscript
pp.~39--41, supplied \texttt{main.tex} lines~3808--4043.
Their complete displayed maps and Poisson-sum argument were
rechecked here. The centering and original operator signs are
those independently derived in MA-MO-007--008.
The two polar characters are genuine finite-dimensional
representation data. In the cote~1 framework, scans~079--080
and~111--112 through ALPOGE\_BRIDGES BR-TR-01/02, the
equivariant exact sequences belong to the additive
representation bookkeeping; a trace or determinant still
requires a defined endomorphism and its analytic trace ideal.
No such trace is assigned to the two-space map (MA-MO.254).
